\section{Введение}
Каскадная обработка позволяет выделять дополнительное вычисление запросам,
на которых оно полезно. Первая модель отвечает всегда, после чего
маршрутизатор сохраняет её ответ либо передаёт запрос второму эксперту.
Поскольку второй эксперт также ошибается, более высокая средняя accuracy
не гарантирует полезность каждого перехода. Цель исследования ---
разработать экспериментально обоснованные рекомендации по выбору способа
обучения маршрутизатора для заданных данных, пары моделей и бюджета.

Объект исследования --- двухмодельный каскад с фиксированными
классификаторами и обучаемым правилом эскалации. На обучающих примерах
известна правильность обеих моделей; при эксплуатации роутер использует
только запрос и доступный результат первой модели. Это позволяет
непосредственно наблюдать спасения ошибок и вредные замены и сравнивать
разные способы оценивания ожидаемой пользы одного и того же перехода.

Сравнительная полезность эскалации формализована в литературе.
\citet{jitkrittum2023confidence} рассматривают разность вероятностей
правильности и прямую регрессию signed gain, а также условия достаточности
confidence-правила. Post-hoc обучение связано с передачей решения
ошибающемуся эксперту \citep{narasimhan2022posthoc}; для языковых моделей
изучены токенные признаки \citep{gupta2024cascades} и выбор продолжения
по ожидаемой полезности \citep{aggarwal2024automix}. Эти работы задают
целевую величину и сильные альтернативные методы для сравнения.

Остаётся практический вопрос выбора оценивателя при конечном объёме
парных данных. Источник и состав обучения по-разному влияют на роутеры
\citep{ong2025routellm,niu2026routing}. В uplift modeling показаны
искажение оценки эффекта после пересэмплирования и роль редких событий
\citep{nyberg2023exploring}. Таким образом, сравнение должно учитывать
одновременно поддержку rescue/harm, смещение отбора, доступную информацию
и квоту, а не только средний разрыв точности экспертов.

Предлагается фиксировать экспертов и признаки, изменять объём и состав
именно обучения роутера и оценивать методы при одинаковой квоте.
Задачи исследования --- выявить режимы предпочтительности методов,
объяснить различия и проверить переносимость рекомендаций на новые
задачи и пары моделей. Дополнительно к accuracy рассматриваются
исправленные ошибки, вредные замены, упущенные спасения и предотвращённый
вред: близкая итоговая точность может скрывать разные профили решений.

Эксперименты проводятся на BANKING77 с предвычисленными прогнозами
отдельно обученных классификаторов. После базовых контролей исследован
полный набор восьми способов обучения на всех 21 прежней паре.
Гипотеза состоит в том, что разные профили спасений и вреда при близкой
accuracy могут сопровождаться неодинаковой чувствительностью к составу
train. Поэтому варьируются направление и сила дисбаланса, а проверяется
изменение \emph{попарной разности} качества методов. Коррекция с известными
вероятностями извлечений служит дополнительным контролем механизма.
Отдельная валидационная часть выбирает метод после настройки
его регуляризации, позволяя оценить и саму процедуру выбора.

Вклад работы --- воспроизводимая проверка условий сравнительного
качества маршрутизаторов: объёма разметки, состава парных исходов
и ошибки выбора по конечному dev. Различаются чувствительность
между семействами целей, различия внутри семи gain-параметризаций
и надёжная смена предпочтительного метода. Сильная сторона протокола
--- фиксированные эксперты, парные исходы и раздельные tuning/selection
данные. Ограничения --- один датасет, зависимые пары, повторное
использование ранее просмотренного test и отсутствие независимой
проверки переносимости рекомендаций.
